\documentclass[a4paper,12pt]{article}
\usepackage[T2A]{fontenc}
\usepackage[utf8]{inputenc}
\usepackage[russian]{babel}
\usepackage{amsmath,amssymb,mathtools}
\usepackage{helvet}
\renewcommand{\familydefault}{\sfdefault}
\usepackage{icomma}
\usepackage[a4paper,margin=20mm,top=22mm,bottom=22mm]{geometry}
\usepackage{microtype}
\usepackage{tikz}
\usetikzlibrary{calc}
\usepackage{tkz-euclide}
\usepackage{tabularx,booktabs}
\usepackage{enumitem}
\usepackage[hidelinks]{hyperref}
\usepackage{parskip}
\usepackage{fancyhdr}
\usepackage{setspace}
\usepackage{xcolor}

\definecolor{accent}{HTML}{008080}
\definecolor{darkaccent}{HTML}{004D4D}
\definecolor{textgray}{HTML}{333333}
\definecolor{softgray}{HTML}{E9EEEE}

\onehalfspacing
\setlength{\parindent}{0pt}
\setlength{\headheight}{25pt}
\setlist[itemize]{leftmargin=1.6em,itemsep=.25em,topsep=.25em}
\setlist[enumerate]{leftmargin=1.8em,itemsep=.45em,topsep=.3em}
\renewcommand{\arraystretch}{1.25}

\pagestyle{fancy}
\fancyhf{}
\fancyhead[L]{\small\color{textgray}Планиметрия: четырёхугольники и прямоугольные треугольники}
\fancyhead[R]{\small\color{textgray}София Кальней}
\fancyfoot[C]{\small\thepage}
\renewcommand{\headrulewidth}{.4pt}

\newcommand{\key}[1]{\textbf{\color{darkaccent}#1}}
\newcommand{\formula}[1]{\[\boxed{#1}\]}

\newcommand{\DrawRightTriangleAltitude}{%
\begin{center}\begin{tikzpicture}
\tkzSetUpLine[line width=.95pt,color=accent]
\tkzSetUpPoint[color=darkaccent,fill=darkaccent,size=3]
\tkzSetUpLabel[color=black,font=\small]
\tkzDefPoint(0,0){A}\tkzDefPoint(5,0){B}\tkzDefPoint(0,3.75){C}\tkzDefPoint(1.8,2.4){H}
\tkzDrawSegment(A,B)\tkzDrawSegment(A,C)\tkzDrawSegment(B,C)\tkzDrawSegment(A,H)
\tkzDrawPoints(A,B,C,H)
\tkzMarkRightAngle[size=.22](B,A,C)\tkzMarkRightAngle[size=.20](A,H,B)
\tkzLabelPoint[below left](A){$A$}\tkzLabelPoint[below right](B){$B$}\tkzLabelPoint[above left](C){$C$}\tkzLabelPoint[above right](H){$H$}
\tkzLabelSegment[below=3pt](A,B){$a$}\tkzLabelSegment[left=3pt](A,C){$b$}\tkzLabelSegment[above right=2pt](B,C){$c$}\tkzLabelSegment[below right=2pt](A,H){$h$}
\end{tikzpicture}\end{center}%
}

\newcommand{\DrawSquareTriangle}[6]{%
\begin{center}\begin{tikzpicture}
\pgfmathsetmacro{\total}{#1+#2+#3}\pgfmathsetmacro{\sfac}{10/\total}\pgfmathsetmacro{\aS}{#1*\sfac}\pgfmathsetmacro{\xS}{#2*\sfac}\pgfmathsetmacro{\nS}{(#1+#2)*\sfac}
\tkzSetUpLine[line width=.95pt,color=accent]\tkzSetUpPoint[color=darkaccent,fill=darkaccent,size=3]\tkzSetUpLabel[color=black,font=\small]
\tkzDefPoint(0,0){A}\tkzDefPoint(10,0){B}\tkzDefPoint(\aS,0){M}\tkzDefPoint(\nS,0){N}\tkzDefPoint(\aS,\xS){L}\tkzDefPoint(\nS,\xS){K}
\tkzInterLL(A,L)(B,K)\tkzGetPoint{C}
\tkzDrawSegment(A,B)\tkzDrawSegment(A,C)\tkzDrawSegment(B,C)\tkzDrawSegment(M,L)\tkzDrawSegment(L,K)\tkzDrawSegment(K,N)\tkzDrawSegment(N,M)
\tkzDrawPoints(A,B,C,M,N,L,K)\tkzMarkRightAngle[size=.22](A,C,B)\tkzMarkRightAngle[size=.17](A,M,L)
\tkzLabelPoint[below left](A){$A$}\tkzLabelPoint[below](M){$M$}\tkzLabelPoint[below](N){$N$}\tkzLabelPoint[below right](B){$B$}\tkzLabelPoint[above](C){$C$}\tkzLabelPoint[above left](L){$L$}\tkzLabelPoint[above right](K){$K$}
\tkzLabelSegment[below=3pt](A,M){#4}\tkzLabelSegment[below=3pt](N,B){#5}
\node[font=\small,text=textgray] at ($(M)!0.5!(K)$) {#6};
\end{tikzpicture}\end{center}%
}

\newcommand{\DrawParallelogram}[4]{%
\begin{center}\begin{tikzpicture}
\pgfmathsetmacro{\base}{2.3}\pgfmathsetmacro{\dx}{.6*#1*\base}\pgfmathsetmacro{\dy}{.8*#1*\base}
\tkzSetUpLine[line width=.95pt,color=accent]\tkzSetUpPoint[color=darkaccent,fill=darkaccent,size=3]\tkzSetUpLabel[color=black,font=\small]
\tkzDefPoint(0,0){A}\tkzDefPoint(\base,0){B}\tkzDefPoint(\dx,\dy){D}\tkzDefPoint({\base+\dx},\dy){C}\tkzDefPoint(1.6,.8){RA}\tkzDefPoint({\base-.4},.8){RB}
\tkzInterLL(A,RA)(D,C)\tkzGetPoint{M}\tkzInterLL(B,RB)(D,C)\tkzGetPoint{N}
\tkzDrawSegment[color=textgray,dashed,line width=.7pt](N,M)
\tkzDrawSegment(A,B)\tkzDrawSegment(B,C)\tkzDrawSegment(C,D)\tkzDrawSegment(D,A)\tkzDrawSegment(A,M)\tkzDrawSegment(B,N)
\tkzDrawPoints(A,B,C,D,M,N)
\tkzMarkAngle[size=.45](B,A,M)\tkzMarkAngle[size=.59](M,A,D)\tkzMarkAngle[size=.45](C,B,N)\tkzMarkAngle[size=.59](N,B,A)
\tkzLabelPoint[below left](A){$A$}\tkzLabelPoint[below right](B){$B$}\tkzLabelPoint[above](N){$N$}\tkzLabelPoint[above](D){$D$}\tkzLabelPoint[above](C){$C$}\tkzLabelPoint[above](M){$M$}
\tkzLabelSegment[below=3pt](A,B){#2}\tkzLabelSegment[below right=2pt](B,C){#3}\tkzLabelSegment[above=4pt](N,M){#4}
\end{tikzpicture}\end{center}%
}

\newcommand{\DrawRhombus}[5]{%
\begin{center}\begin{tikzpicture}
\pgfmathsetmacro{\side}{4.5}\pgfmathsetmacro{\dx}{\side*cos(#1)}\pgfmathsetmacro{\dy}{\side*sin(#1)}\pgfmathsetmacro{\ox}{(\side+\dx)/2}\pgfmathsetmacro{\oy}{\dy/2}
\tkzSetUpLine[line width=.95pt,color=accent]\tkzSetUpPoint[color=darkaccent,fill=darkaccent,size=3]\tkzSetUpLabel[color=black,font=\small]
\tkzDefPoint(0,0){A}\tkzDefPoint(\side,0){B}\tkzDefPoint(\dx,\dy){D}\tkzDefPoint({\side+\dx},\dy){C}\tkzDefPoint(\ox,\oy){O}\tkzDefPoint(\ox,0){H}
\tkzDrawSegment(A,B)\tkzDrawSegment(B,C)\tkzDrawSegment(C,D)\tkzDrawSegment(D,A)\tkzDrawSegment(A,C)\tkzDrawSegment(B,D)\tkzDrawSegment(O,H)
\tkzDrawPoints(A,B,C,D,O,H)\tkzMarkRightAngle[size=.18](O,H,B)\tkzMarkRightAngle[size=.16](A,O,B)\tkzMarkAngle[size=.62](B,A,D)
\tkzLabelPoint[below left](A){$A$}\tkzLabelPoint[below right](B){$B$}\tkzLabelPoint[above left](D){$D$}\tkzLabelPoint[above right](C){$C$}\tkzLabelPoint[above right](O){$O$}\tkzLabelPoint[below left](H){$H$}
\tkzLabelSegment[below=3pt](A,B){#2}\tkzLabelSegment[left=3pt](O,H){#5}
\if\relax\detokenize{#4}\relax\else\tkzLabelAngle[pos=.92](B,A,D){#4}\fi
\if\relax\detokenize{#3}\relax\else\node[font=\small,text=textgray] at ($(O)+(0,.62)$) {#3};\fi
\end{tikzpicture}\end{center}%
}

\begin{document}
\hypersetup{pageanchor=false}
\begin{titlepage}\thispagestyle{empty}\centering
\vspace*{2.4cm}
{\Large\color{darkaccent}\textbf{Чек-лист}}\par\vspace{.7cm}
{\huge\bfseries Планиметрия повышенной сложности}\par\vspace{.35cm}
{\Large Прямоугольный треугольник, параллелограмм, ромб}\par\vspace{1.7cm}
{\large Подобие · тригонометрия · биссектрисы · площади · дополнительные построения}\par\vspace{2.2cm}
{\large\textbf{Лёвин Артём Александрович}}\par\vspace{.25cm}
{\large эксклюзивно для Софии Кальней}\par\vspace{.8cm}
{\large 4 сентября 2026 года}\par\vfill
\begin{tikzpicture}[x=1cm,y=1cm]
\draw[accent,line width=1.1pt] (0,0) .. controls (2.2,.65) and (4.6,-.55) .. (6.8,.05) .. controls (8.7,.55) and (10.6,.35) .. (12.8,.75);
\draw[darkaccent,line width=.45pt] (.7,-.25) .. controls (3.1,.15) and (5.2,-.35) .. (7.2,-.05) .. controls (9.2,.25) and (10.9,.05) .. (12.2,.35);
\end{tikzpicture}\end{titlepage}
\hypersetup{pageanchor=true}

\section*{\color{darkaccent}1. Базовый алгоритм сложной планиметрии}
\begin{enumerate}[label=\textbf{\arabic*.}]
\item \key{Найдите треугольники.} В первую очередь прямоугольные: для них доступны теорема Пифагора и тригонометрия.
\item \key{Если треугольников много --- ищите подобие.} Сигналы: вертикальные, накрест лежащие, общие и прямые углы.
\item \key{Если есть параллельные прямые --- ищите угловые связи.} В параллелограмме это один из главных ресурсов.
\item \key{Если есть биссектриса --- ищите равнобедренный треугольник.} Биссектриса вместе с параллельностью часто создаёт пару равных углов.
\item \key{Если нужно расстояние от точки до прямой --- стройте перпендикуляр.}
\item \key{Если решение быстро усложняется, смените метод.} Появление нескольких лишних неизвестных и корней --- сигнал проверить подобие, тригонометрию или площади.
\end{enumerate}

\subsection*{Прямоугольный треугольник}
Проверяйте три инструмента: теорему Пифагора, тригонометрические отношения и специальные формулы.
\DrawRightTriangleAltitude
Если катеты равны $a,b$, гипотенуза $c$, высота к гипотенузе $h$, то
\[\frac{ab}{2}=\frac{ch}{2},\qquad \boxed{h=\frac{ab}{c}}.\]

\section*{\color{darkaccent}2. Квадрат в прямоугольном треугольнике}
Пусть $\angle C=90^\circ$, $M,N\in AB$, $L\in AC$, $K\in BC$, а $KLMN$ --- квадрат. Обозначим
\[AM=a,\qquad BN=b,\qquad MN=KN=LM=KL=x.\]
\DrawSquareTriangle{1}{2}{4}{$a$}{$b$}{$x$}
\subsection*{Метод 1. Подобие}
Треугольники $KNB$ и $MAL$ прямоугольные, причём $\angle KBN=\angle MLA$. Следовательно,
\[\triangle KNB\sim\triangle MAL.\]
По сходственным сторонам
\[\frac{KN}{AM}=\frac{BN}{LM}.\]
Подставляя $KN=LM=x$, получаем
\[\frac{x}{a}=\frac{b}{x},\qquad \boxed{x^2=ab}.\]
Поскольку площадь квадрата равна $x^2$,
\formula{S_{\text{кв}}=ab}
\subsection*{Метод 2. Через тангенс}
Из равенства острых углов:
\[\tg\angle KBN=\frac{x}{b},\qquad \tg\angle MLA=\frac{a}{x}.\]
Поэтому $x/b=a/x$, откуда $x^2=ab$.
\key{Вывод.} Если несколько прямоугольных треугольников имеют одинаковый острый угол, проверяйте подобие или равенство тангенсов.

\section*{\color{darkaccent}3. Параллелограмм и биссектрисы}
Пусть $ABCD$ --- параллелограмм, $BC=AD=2AB$. Биссектрисы углов $A$ и $B$ пересекают \emph{прямую} $CD$ в точках $M,N$.
\DrawParallelogram{2}{$x$}{$2x$}{$MN$}
Положим $AB=x$, тогда $AD=BC=2x$. Биссектриса вместе с параллельностью даёт равные углы, поэтому отсекает равнобедренный треугольник. В результате
\[DN=x,\qquad DC=x,\qquad CM=x,\]
и
\formula{MN=3x}
Если $MN=12$, то $x=4$, поэтому
\[\boxed{AB=CD=4,\qquad BC=AD=8}.\]
\subsection*{Обобщение из занятия}
Если $BC=k\,AB$, $k>1$, и $AB=x$, то
\[DN=(k-1)x,\qquad DC=x,\qquad CM=(k-1)x,\]
следовательно,
\formula{MN=(2k-1)x}
\key{Важно.} «Прямая $CD$» продолжается в обе стороны, поэтому $M,N$ могут лежать за пределами стороны $CD$.

\section*{\color{darkaccent}4. Ромб: расстояние от центра до стороны}
Пусть сторона ромба равна $a$, острый угол равен $\alpha$, $O$ --- точка пересечения диагоналей, $OH$ --- расстояние от центра до стороны.
\DrawRhombus{30}{$a$}{}{$\alpha$}{$OH$}
По определению $OH\perp AB$. Высота ромба
\[h=a\sin\alpha.\]
Центр ромба находится посередине между двумя параллельными сторонами, поэтому
\formula{OH=\frac{a\sin\alpha}{2}}
Для $a=4$, $\alpha=30^\circ$ получаем $OH=1$.
\subsection*{Проверка через площадь}
\[S=a^2\sin\alpha=ah.\]
Отсюда $h=a\sin\alpha$, затем $OH=h/2$.
\subsection*{Ещё один путь}
При $\angle A=30^\circ$ диагональ даёт $\angle BAO=15^\circ$. В прямоугольном $\triangle AOB$:
\[BO=AB\sin15^\circ,\qquad AO=AB\cos15^\circ.\]
Высота к гипотенузе:
\[OH=\frac{AO\cdot BO}{AB}=AB\sin15^\circ\cos15^\circ=\frac{AB\sin30^\circ}{2}.\]

\section*{\color{darkaccent}5. Формулы и признаки из памяти}
\begin{itemize}
\item $a^2+b^2=c^2$ --- теорема Пифагора.
\item $h=ab/c$ --- высота к гипотенузе.
\item $\sin\alpha=$ противолежащий катет / гипотенуза; $\cos\alpha=$ прилежащий катет / гипотенуза; $\tg\alpha=$ противолежащий / прилежащий катет.
\item Вертикальные углы равны; при параллельных прямых накрест лежащие углы равны.
\item В параллелограмме $AB=CD$, $BC=AD$, $AB\parallel CD$, $BC\parallel AD$.
\item В ромбе все стороны равны; диагонали перпендикулярны, делят друг друга пополам и делят углы пополам.
\item $S=ah$, а для ромба также $S=a^2\sin\alpha$.
\item Расстояние от точки до прямой --- длина перпендикуляра к прямой.
\end{itemize}

\section*{\color{darkaccent}6. Сигналы на чертеже}
\begin{tabularx}{\linewidth}{@{}p{.35\linewidth}X@{}}
\toprule
\textbf{Что увидели} & \textbf{О чём подумать}\\
\midrule
Прямоугольный треугольник & Пифагор, $\sin$, $\cos$, $\tg$, высота к гипотенузе\\
Много треугольников & Подобие или равенство\\
Вертикальные углы & Возможное подобие по двум углам\\
Параллельные прямые & Накрест лежащие и соответственные углы\\
Биссектриса в параллелограмме & Возможный равнобедренный треугольник\\
Неизвестная высота & Площадь одной фигуры двумя способами\\
Нужно расстояние от точки до прямой & Провести перпендикуляр\\
Решение породило много неизвестных & Проверить другой метод\\
\bottomrule
\end{tabularx}

\section*{\color{darkaccent}7. Типичные ошибки}
\begin{itemize}
\item Путать сторону $CD$ и прямую $CD$.
\item Утверждать подобие без указания двух пар равных углов.
\item Сопоставлять сходственные стороны только по положению на рисунке.
\item Считать любой отрезок от точки до прямой расстоянием; требуется перпендикуляр.
\item Не использовать параллельность сторон параллелограмма для получения равных углов.
\item Забывать, что диагональ ромба делит угол пополам.
\item Продолжать громоздкий Пифагор, когда подобие, тангенс или площади дают более короткий путь.
\end{itemize}

\newpage
\section*{\color{darkaccent}Тренировочный блок}
\textbf{Средний уровень}
\begin{enumerate}[label=\textbf{\arabic*.}]
\item В прямоугольном треугольнике $ABC$ квадрат $KLMN$ расположен как в теории. $AM=3$, $BN=12$. Найдите площадь квадрата.
\DrawSquareTriangle{3}{6}{12}{$3$}{$12$}{$S_{\text{кв}}=?$}
\item В той же конфигурации $AM=5$, $BN=20$. Найдите сторону квадрата и его площадь.
\DrawSquareTriangle{5}{10}{20}{$5$}{$20$}{$x=?,\ S=?$}
\item Сторона ромба равна $10$, острый угол $30^\circ$. Найдите расстояние от центра ромба до стороны.
\DrawRhombus{30}{$10$}{}{$30^\circ$}{$?$}
\end{enumerate}
\textbf{Выше среднего}
\begin{enumerate}[label=\textbf{\arabic*.},start=4]
\item В параллелограмме $BC=2AB$. Биссектрисы углов $A,B$ пересекают прямую $CD$ в $M,N$, $MN=21$. Найдите стороны параллелограмма.
\DrawParallelogram{2}{$AB=?$}{$BC=2AB$}{$MN=21$}
\item Сторона ромба $8$, острый угол $60^\circ$. Найдите расстояние от центра ромба до стороны.
\DrawRhombus{60}{$8$}{}{$60^\circ$}{$?$}
\item В конфигурации с квадратом $S_{\text{кв}}=54$, $AM=6$. Найдите $BN$.
\DrawSquareTriangle{6}{7.348469228}{9}{$6$}{$?$}{$S_{\text{кв}}=54$}
\item В параллелограмме $BC=3AB$, $MN=25$. Найдите $AB$ и $BC$.
\DrawParallelogram{3}{$AB=?$}{$BC=3AB$}{$MN=25$}
\item Площадь ромба $24$, сторона $6$. Найдите расстояние от центра ромба до стороны.
\DrawRhombus{41.810315}{$6$}{$S=24$}{}{$?$}
\item В конфигурации с квадратом $AM=2\sqrt3$, $BN=6\sqrt3$. Найдите сторону квадрата.
\DrawSquareTriangle{3.464101615}{6}{10.392304845}{$2\sqrt3$}{$6\sqrt3$}{$x=?$}
\end{enumerate}
\textbf{Сложный уровень}
\begin{enumerate}[label=\textbf{\arabic*.},start=10]
\item В параллелограмме $BC=3AB$, $MN=35$. Найдите периметр параллелограмма.
\DrawParallelogram{3}{$AB=?$}{$BC=3AB$}{$MN=35$}
\end{enumerate}

\newpage
\section*{\color{darkaccent}Ответы}
\begin{table}[h!]\centering
\caption{Ответы к тренировочному блоку}
\begin{tabularx}{\linewidth}{@{}cX@{}}
\toprule
\textbf{№} & \textbf{Ответ}\\
\midrule
1 & $36$\\
2 & сторона $10$, площадь $100$\\
3 & $\dfrac52$\\
4 & $AB=CD=7$, $BC=AD=14$\\
5 & $2\sqrt3$\\
6 & $BN=9$\\
7 & $AB=5$, $BC=15$\\
8 & $2$\\
9 & $6$\\
10 & $56$\\
\bottomrule
\end{tabularx}
\end{table}
\end{document}
